\documentclass[conference]{IEEEtran}
\IEEEoverridecommandlockouts
\usepackage{cite}
\usepackage{graphicx}
\usepackage{booktabs}
\usepackage{url}
\begin{document}
\title{A Software-Defined Low-Cost CAN Testbed for Automotive Cybersecurity Education and Forensic-Readiness Analysis}
\author{\IEEEauthorblockN{Osvaldo Janeri Filho}\IEEEauthorblockA{Independent Researcher / Digital Forensics Practitioner\\Fortaleza, Cear\'a, Brazil\\Email omitted for blind or camera-ready submission}}
\maketitle
\begin{abstract}
Modern vehicles rely on distributed electronic control units connected through in-vehicle networks such as the Controller Area Network (CAN), whose legacy design lacks native authentication and message-origin verification. This paper presents a software-defined, low-cost CAN testbed for automotive cybersecurity education and forensic-readiness analysis. The artifact combines virtual CAN traffic generation, controlled spoofing and flooding scenarios, normalized datasets, evidence manifests, SHA-256 hash preservation, and a client-side observer demo. Two software-only experiments generated a combined dataset of 42,288 labeled CAN frames. A rule-based triage heuristic is used only as a label-consistency and forensic-screening sanity check, not as a real-vehicle intrusion detection benchmark.
\end{abstract}
\begin{IEEEkeywords}
automotive cybersecurity, CAN bus, forensic readiness, testbed, virtual CAN, reproducibility
\end{IEEEkeywords}
% This is a starter IEEEtran skeleton. Port the content from article_ieee_v2.html section by section before final submission.
\section{Introduction}
TODO: Port revised HTML v2 content.
\section{Background and Related Work}
TODO.
\section{Testbed and Evidence Workflow}
TODO.
\section{Experimental Methodology}
TODO.
\section{Results}
TODO.
\section{Triage Sanity Check}
TODO.
\section{Discussion and Limitations}
TODO.
\section{Conclusion}
TODO.
\bibliographystyle{IEEEtran}
\bibliography{references}
\end{document}
